\section{Introduction}
\label{sec:intro}

REST APIs carry most of the traffic between modern software components, and testing them
well is expensive and still largely manual. LLMs that read an OpenAPI specification and emit
executable tests are therefore of obvious commercial interest, and the research response has
been quick. RESTGPT~\cite{restgpt} extracts input-constraint rules at 97\% precision.
AutoRestTest~\cite{autoresttest} reports 42 server errors against EvoMaster's 20,
LlamaRestTest~\cite{llamaresttest} reports 204 faults against 130, and
RESTifAI~\cite{restifai} reaches 95.5\% operation coverage.

These numbers answer a narrower question than they appear to. Each counts what a suite found
on a live system, and none expresses that count as a fraction of what was available to find.
Without a known fault population, ``42 versus 20'' cannot distinguish a suite that found 42
of 50 faults from one that found 42 of 500. For a QA lead deciding whether to trust a
generator, the missing denominator is the whole question.

\subsection{Gaps}
Our systematic review of 59 papers (2018--2026, five independent searches merged into
\path{team-synthesis/evidence-table-merged.md}) isolates three persistent gaps.

\textbf{GAP-C (Comparison).} Only one study compares a generator against a manual suite:
EvoMaster~\cite{evomaster}, which found its own tool below manual (41\% vs 82\% statement
coverage). That study predates LLMs. No paper places an LLM, a manual suite and a mature tool
on the same APIs, so the field's central practical claim, that LLM tests rival human ones,
has never been tested directly.

\textbf{GAP-D (Dataset / ground truth).} Fault detection is universally reported as raw
500-error counts on live
APIs~\cite{llamaresttest,autoresttest,logiagent,notimetorest,evomaster,morest}. With no
shared pre-seeded-fault benchmark, ground-truth Recall is never computed.

\textbf{GAP-M (Metric).} Evaluation is coverage-biased. Only KAT~\cite{kat} splits 2xx from
4xx and only RESTSpecIT~\cite{restspecit} reports 5xx. No work reports a per-endpoint
edge-case metric or an endpoint-type miss profile.

\subsection{This Paper}
We close GAP-C and GAP-D together. Seeding faults with standard mutation operators into three
EMB APIs fixes the denominator at \RQIIn{} known faults, and we then run all three arms
against every mutant. For GAP-M we add a per-endpoint edge-case count and a code-derived
\emph{error-surface baseline}: a static enumeration, per endpoint, of every source location
able to emit a non-2xx response. This gives a tool-independent map of what can fail, against
which each arm's behaviour is read.

The outcome divides along an unexpected line. On the metrics the literature favours, the LLM
wins by a wide margin, with \RQIcoveragePct\% endpoint coverage and \RQIIItotalLlm{} edge-case
scenarios against the manual suite's \RQIIItotalManual{} ($p=\RQIIIpvalue$). On the ground
truth the ranking is EvoMaster \RQIIrecallEvo{} $>$ LLM \RQIIrecallLlm{} $>$ manual
\RQIIrecallManual{}. The LLM beats an independently authored human-style suite and is beaten
by an established search-based tool, writing roughly \RQIIItotalLlm{} scenarios while
detecting fewer faults than that tool.

Section~\ref{sec:discussion} traces this dissociation to two mechanisms that compose. Oracle
\emph{kind} sets a ceiling, because the spec-derived status-code assertions written by the LLM
and by the manual suite cannot see a fault that changes a computed value while leaving the
status untouched, which is what most of our mutants do. Within that status-oracle regime,
volume decides the rest. Together the two levels explain how a suite can be more thorough by
every published metric, better than a human-style baseline, and still less effective than a
mature tool.

\subsection{Contributions}
\begin{enumerate}
  \item \textbf{The first three-way comparison} of LLM, manual and EvoMaster test generation
        on identical APIs under a blind protocol (GAP-C).
  \item \textbf{A pre-seeded-fault benchmark} of \RQIIn{} catalogued mutants over
        \RQIopsTotal{} operations, enabling true fault-detection Recall rather than live-API
        error counts (GAP-D).
  \item \textbf{A per-endpoint edge-case metric and a code-derived error-surface baseline},
        which separate scenarios produced from faults detected (GAP-M).
  \item \textbf{A two-level mechanism} in which oracle kind sets the ceiling and oracle volume
        allocates what falls under it, explaining why volume and effectiveness diverge and why
        scenario-count metrics mislead when reported alone.
  \item \textbf{A reproducible pipeline} in which every number in this paper is generated from
        committed raw data by a script rather than transcribed
        (Section~\ref{sec:method-repro}).
\end{enumerate}
